% FILE: 01_research_question.tex
% Project: lifecycle
% Role: process-lifecycle-state-definitions-and-transitions

\section{Research Question}
\label{sec:lifecycle-question}

\subsection{Problem Statement}
\label{sec:lifecycle-question-problem}

Enterprise process intelligence, as practiced in M\&A due diligence and operational
governance, rests on claims that are structurally unverifiable. A seller states that
their order-to-cash cycle time is 3.2 days, their compliance rate is 98.7\%, and their
process is ISO~9001 certified. The buyer reviews a deck, interviews management, and
accepts or rejects these claims based on documentation and interviews. The event log that
produced the ``3.2 days'' figure is not provided. The ``98.7\% compliance rate'' cannot
be reproduced because no conformance check is available. The ISO certification refers to
a point-in-time audit, not continuous conformance. The process model used in the
certification may not match the process actually running in production.

This is the diligence gap, defined precisely in \texttt{lifecycle/ACQUISITION.md}: the
claims are real to the seller in the sense that someone computed them, but they are not
replayable by the buyer. The buyer cannot distinguish an accurate claim from an
aspirational one. In the absence of replayable evidence, the buyer applies a risk premium
to cover unquantified uncertainty. The risk premium is not punitive; it reflects real
epistemic debt. The research problem addressed by the lifecycle project is: what formal
state machine, with what admissibility criteria, would make process claims replayable by
construction rather than by negotiation?

\subsection{Old-Regime Assumption Challenged}
\label{sec:lifecycle-question-old-regime}

The old-regime assumption challenged by the lifecycle framework is that a process model
is a document---a BPMN diagram, a flowchart, a slide---that describes intended behavior
and is updated by human editors when behavior changes. Under this assumption, process
governance is a documentation discipline: maintaining correspondence between the diagram
and the process. Conformance is audited periodically, manually, and reported as a
narrative: ``we believe the process works as follows.''

The lifecycle framework rejects this assumption. A process model is not a document; it is
a typed artifact with a formally specified state. A WF-net must carry a
\texttt{WfNetConst<SOUNDNESS>} type witness. An event log must be admitted through
\texttt{Admit::admit()} before it can participate in conformance computation. A
conformance score must be expressed as \texttt{Metric<Fitness, NUM, DEN>} bounded by
\texttt{Between01<NUM, DEN>}. A repair action must produce a \texttt{Receipt}. Without
these type constraints, the process model is L1 Ad-hoc by definition: no formal models,
no systematic event log, no replayable conformance. Governance claims are narration.

The old regime further assumes that autonomic behavior is an advanced feature to be added
after basic governance is established. The lifecycle framework challenges this assumption
by treating autonomic closure---the full MAPE-K loop with receipted artifacts at every
stage---as the standard that all formal process governance must eventually reach, with
maturity levels L1 through L5 measuring progress toward that standard.

\subsection{Hypothesis}
\label{sec:lifecycle-question-hypothesis}

The dissertation hypothesis, as instantiated in the lifecycle project, is as follows:
a formally specified 13-state lifecycle, governed by six quality gates derived from
Petri net soundness theory and alignment-based conformance checking, is sufficient to
make process intelligence claims replayable, board-defensible, and autonomically
maintainable without requiring human intervention at each conformance cycle.

More precisely, the hypothesis is that the combination of (1) a typed evidence lifecycle
$\mathcal{S} = \{\texttt{Raw}, \ldots, \texttt{Receipted}\}$ with type-level state
enforcement, (2) six gate criteria formalized as mathematical predicates over WF-nets and
event logs, (3) a MAPE-K loop mapping that assigns each operational stage to its autonomic
function, and (4) a cryptographic receipt chain that accumulates evidence across the full
lifecycle, produces a system in which every governance claim resolves to a verifiable
artifact. The hypothesis is stated as PARTIAL in the lifecycle project's current
checkpoint status: the doctrine is complete and the type definitions exist in
\texttt{wasm4pm-compat}, but the execution engine (\texttt{wasm4pm}) is a separate
downstream product that has not yet closed the loop at runtime.

\subsection{Success Criteria}
\label{sec:lifecycle-question-success}

The success criteria for the lifecycle project are defined in
\texttt{lifecycle/checkpoint\_\_lifecycle\_model\_complete.md} as a conjunction of four
automated assertions. First, \texttt{assert\_soundness(N)}: the process model $N$ is a
valid WF-net, 1-bounded, live, and free of dead transitions. Second,
\texttt{assert\_fitness(N, L, 0.95)}: the alignment fitness $\mathrm{fitness}(L, N) \geq
0.95$ against the active event log. Third,
\texttt{assert\_no\_ghost\_transitions(N, L)}: every labeled transition in the net is
fired at least once in the log or is classified as an explicit silent ($\tau$) transition.
Fourth, \texttt{assert\_decommission\_receipt(R\_d, K\_\mathrm{pub})}: the
cryptographic decommissioning receipt $R_d$ is structurally valid and carries a verified
signature against the model hash and final log hash.

The program-level success criterion is the ALIVE verdict issued in
\texttt{checkpoints/PROCESS\_INTELLIGENCE\_ALIVE\_001.md} with verification status code
\texttt{0x00} (SUCCESS\_ALIVE). This checkpoint confirms that the mathematical invariant
checklist---admissibility boundary, autonomic actuation, token game fitness, OCPQ
refinement, and decommissioning retirement---has been validated at the research program
level. The lifecycle-module-specific ALIVE verdict, requiring the four assertions above
to execute against real Petri nets and real event logs in a running \texttt{wasm4pm}
instance, remains a deferred gate.
